\documentclass{article}
\usepackage[UTF8]{ctex}
\usepackage{amsmath,amssymb,amsfonts}
\usepackage{enumitem}
\usepackage{geometry}
\geometry{a4paper, margin=1in}
% === 错题本排版优化（不参与题目标识） ===
\setlength{\parindent}{0pt}
\setlength{\parskip}{0.35em plus 0.10em minus 0.08em}
\setlength{\emergencystretch}{3em}
\setlist{itemsep=0.25em, topsep=0.40em, parsep=0pt}
\allowdisplaybreaks[2]
% === 错题本排版优化结束 ===
\begin{document}
\section*{苏州大学2025年高等代数真题}
\section*{82．苏州大学2025年高等代数试卷}

\subsection{题目1}
1．解答如下问题。

\begin{itemize}
  \item[（1）] 设 \(\mu\) 为实数，记
\[
A=\left(\begin{array}{lll}
1 & 0 & 1 \\
1 & 1 & \mu \\
0 & 0 & 1
\end{array}\right) .
\]
求 \(A\) 的若尔当标准形 \(J\) 。
  \item[（2）] 求 \(\mu\) 为何值时，存在正交矩阵 \(P\) ，使得 \(P^{-1} A P=J\) 。
\end{itemize}

\subsection{题目2}
2．设 \(f(x)\) 为数域 \(F\) 上的 \(n \geq 1\) 次多项式，\(f^{\prime}(x)\) 为 \(f(x)\) 的导数．证明：\(f^{\prime}(x) \mid f(x)\) 当且仅当 \(f(x)\)有 \(n\) 重根。

\subsection{题目3}
3．设 \(V\) 为数域 \(F\) 上的 \(n\) 维线性空间，\(\alpha_{1}, \alpha_{2}, \cdots, \alpha_{n}\) 为 \(V\) 的一组基。证明：\\
（1）对于 \(V\) 的任意 \(m\) 个非平凡子空间 \(V_{1}, V_{2}, \cdots, V_{m}\) ，有 \(\bigcup_{i=1}^{m} V_{i} \neq V\) ．\\
（2）对任意的正整数 \(s\) ，存在 \(\alpha_{n+1}, \cdots, \alpha_{n+s} \in V\) ，使得向量组 \(\alpha_{1}, \cdots, \alpha_{n}, \alpha_{n+1}, \cdots, \alpha_{n+s}\) 中任意 \(n\) 个向量够构成 \(V\) 的一组基．

\subsection{题目4}
4．设 \(V\) 为薮域 \(F\) 上的 \(2 n\) 维线性空间，\(\sigma\) 为 \(V\) 上的线性变换．证明： \(\operatorname{Ker} \sigma=\operatorname{Im} \sigma\) 当且仅当存在 \(V\)的一组基，使得 \(\sigma\) 在这组基下矩阵为 \(\left(\begin{array}{cc}O & E_{n} \\ O & O\end{array}\right)\) ．

\subsection{题目5}
5．设 \(\varepsilon_{1}, \varepsilon_{2}, \cdots, \varepsilon_{n}\) 为有理数域 \(\mathbb{Q}\) 上 \(n\) 维线性空间 \(V\) 的一组基，\(V\) 上线性变换 \(\sigma\) 在这组基下的矩阵为

\[
\left(\begin{array}{cc}
0 & 1 \\
E_{n-1} & 0
\end{array}\right)
\]

令 \(U=\left\{\sum_{i=1}^{n} k_{i} \varepsilon_{i} \mid \sum_{i=1}^{n} k_{i}=0, k_{i} \in \mathbb{Q}\right\}, W=\left\{k \sum_{i=1}^{n} \varepsilon_{i}, k \in \mathbb{Q}\right\}\) 。

\begin{itemize}
  \item[（1）] 证明：\(V=W \oplus U\) 。
  \item[（2）] 当 \(n\) 为素数时，求 \(V\) 的所有 \(\sigma-\) 子空间．
\end{itemize}

\subsection{题目6}
6．设 \(V\) 为 \(n\) 维欧氏空间，对于任意的 \(\alpha \in V, \alpha \neq 0\) ，定义 \(\sigma_{\alpha}(\gamma)=\gamma-2 \frac{(\gamma, \alpha)}{(\alpha, \alpha)} \alpha, \gamma \in V\) ，其中（）表示内积，\(\sigma_{\alpha}\) 称为镜面反射．

\begin{itemize}
  \item[（1）] 证明：\(\sigma_{\alpha}\) 是正交变换．
  \item[（2）] 设 \(\eta_{1}, \eta_{2}, \eta_{1}^{\prime}, \eta_{2}^{\prime}\) 为 \(V\) 中的单位向量，\(\left(\eta_{1}, \eta_{2}\right)=\left(\eta_{1}^{\prime}, \eta_{2}^{\prime}\right)\) ，证明：存在 \(V\) 上的线性变换 \(\tau\) ，使得
  \item[（i）] \(\tau\left(\eta_{i}\right)=\eta_{i}^{\prime}, i=1,2\) ．
  \item[（ii）] \(\tau\) 为镜面反射或者两个镜面反射的乘积．
\end{itemize}

\subsection{题目7}
7．设 \(\sigma, \tau\) 为 \(n\) 维复线性空间 \(V\) 上的线性变换，满足 \(\sigma \tau-\tau \sigma=\tau\) 。证明：

\begin{itemize}
  \item[（1）] 若 \(W\) 为 \(\sigma-\) 子空间，则 \(\tau(W)\) 仍是 \(\sigma-\) 子空间。
  \item[（2）] 存在 \(V\) 的一组基，使得 \(\sigma, \tau\) 在这组基下的矩阵都是下三角矩阵。
\end{itemize}
\end{document}
